\section{Appendix F: Independent Review Findings and Resolutions}\label{sec:appendix-review-resolutions}

Following the revision documented in Appendix~E, the thesis underwent an
independent technical review pass covering the communications theory, the
machine-learning claims, the numerical results, and their internal consistency.
The review raised twenty-four findings, which were recorded as inline
annotations in the working draft. \textbf{Every finding has now been resolved
directly in the text}; this appendix documents each one and the change made, so
that the resolutions are auditable without reference to the annotated draft.

Findings are grouped by type. ``Substantive'' denotes a change to a technical
claim, a derivation, or a stated result; ``consistency'' denotes reconciliation
of numbers or cross-references that disagreed between chapters; ``editorial''
denotes typography, notation, or presentation.

\subsection*{F.1 Substantive corrections}

\begingroup\small
\begin{enumerate}\setlength{\itemsep}{3pt}\setlength{\parskip}{1pt}

\item \textbf{AF end-to-end SNR stated without attribution or assumptions}
(Section~\ref{sec:classical-relay-strategies-af}). Equation~\ref{eq:af-snr} is a
standard result but carried no citation, and the $+1$ term in its denominator is
convention-dependent. \emph{Resolved:} the expression is now attributed, and the
text states explicitly that it assumes CSI-assisted \emph{variable-gain} AF (the
variant implemented here), explains the origin of the $+1$ term, and
distinguishes it from the fixed-gain convention.

\item \textbf{Overclaim on universal approximation}
(Section~\ref{sec:theoretical-basis-universal-approximation-and-denoising}). The
text asserted that a single-hidden-layer network with tanh output ``can represent
this function exactly.'' This holds only for a \emph{fixed} noise variance, whereas
the deployed relay is trained at $\{5,10,15\}$~dB and evaluated over $0$--$20$~dB,
and on fading must cover a mixture of random effective SNRs. \emph{Resolved:} the
claim is now split (exact representation for fixed $\sigma^2$, close
\emph{approximation} of the multi-SNR family) and the latter is identified as an
empirical finding rather than a theoretical consequence.

\item \textbf{Unjustified contraction claim in the Mamba selectivity mechanism}
(Section~\ref{sec:mamba-selective-state-spaces}). The step
$\exp(\Delta_k\mathbf{A})\to\mathbf{0}$ for large $\Delta_k$ requires
$\mathbf{A}$ to be Hurwitz, which was never stated. \emph{Resolved:} the text now
states the parameterization $\mathbf{A}=-\exp(A_{\log})$ explicitly, notes that it
makes $\mathbf{A}\prec 0$ diagonal, and derives the reset/retain behaviour from
that property.

\item \textbf{VAE result attributed to the generative paradigm without ruling out
an inference-time cause} (Sections~\ref{sec:generative-vs.-discriminative-paradigms-for-relay-processing}
and~\ref{sec:siso-bpsk-performance-baseline-relay-comparison}). The VAE relay is
pinned near BER $0.25$--$0.40$ at all SNRs, and this was read as evidence about
generative modelling as such. The relay samples $\mathbf{z}\sim q_\phi$
stochastically at inference; the standard deterministic variant (posterior mean
$\boldsymbol{\mu}$ at test time) was never evaluated. \emph{Resolved:} both the
literature-review conclusion and the results discussion now state this caveat
explicitly and report the VAE outcome as specific to the evaluated configuration,
with the deterministic-decoding comparison identified as future work.

\item \textbf{``Identifiability'' mischaracterized in the pilot-budget sweep}
(Section~\ref{sec:partial-posterior}). The collapse of pilot-aided Viterbi at
5~pilots was described as the channel becoming unidentifiable, but with
5~observations and 3~complex taps the least-squares problem remains
overdetermined. \emph{Resolved:} the mechanism is now correctly described as an
\emph{estimation-variance} cliff (too few observations to average the noise, so
occasional catastrophic fits dominate) and the confidence interval at that point
($0.119\pm0.084$, versus $\pm0.002$ at 10~pilots) is reported in both the text and
the figure caption as direct evidence.

\item \textbf{Uniqueness overreach in the block-length discussion}
(Section~\ref{sec:partial-posterior}). The learned relay was called ``the only
option that remains reliable'' on short blocks, but blind CMA is also pilot-free
and was not tested at $L=40$. \emph{Resolved:} the claim is narrowed to the only
\emph{one-shot} option, the distinguishing property (zero per-block adaptation) is
stated, and the untested CMA behaviour at short block length is flagged as open
rather than asserted.

\item \textbf{Wall-clock speed-up presented without implementation context}
(Section~\ref{sec:complexity-viterbi-mlp}). The $30$--$90\times$ figure compares a
vectorized NumPy matrix multiply against an interpreted Python trellis loop, so
part of the ratio is interpreter overhead. \emph{Resolved:} the figure is now
explicitly scoped to the reference implementations, notes that an optimized C or
SIMD Viterbi would narrow the gap, and defers the load-bearing conclusion to the
implementation-independent $M^{L}$-versus-constant scaling argument.

\item \textbf{Scope statement contradicted the principal contribution}
(Section~\ref{sec:scope-and-delimitations}). The delimitation ``perfect CSI is
assumed; channel estimation errors are not modeled'' was stale, since
Chapter~\ref{sec:unknown-channels} models finite-pilot LS estimation and a
zero-CSI blind regime. \emph{Resolved:} the item is now scoped to the canonical
core, and the lifted assumptions in the unknown-channel chapter are named.

\item \textbf{Displayed system model showed the calibration case, not the
canonical channel} (Section~\ref{sec:relay-processing-methods}).
Equation~\ref{eq:relay-received-siso} showed a fixed noise variance although the
canonical channel is Rayleigh, giving a random per-symbol effective SNR.
\emph{Resolved:} the equation is rewritten in post-compensation form
$\tilde{y}_R[i]=|h_1[i]|x[i]+\tilde{n}_1[i]$, the exponentially distributed
effective SNR $|h_1|^2\gamma$ is stated, and the AWGN limit is identified as the
$|h_1|\equiv 1$ special case.

\item \textbf{Inexact DF error-composition formula}
(Section~\ref{sec:classical-relay-strategies-df}). The chapter gave
$P_{e,1}+(1-P_{e,1})P_{e,2}$ in one place and the exact
$P_{e,1}+P_{e,2}-2P_{e,1}P_{e,2}$ in another. \emph{Resolved:} both now use the
exact odd-flip form, with a sentence explaining that two errors cancel, so the two
appearances agree by construction.

\item \textbf{MIMO equalization section implied unreported experiments.} Chapter~2
carried a ``Theoretical Foundations: MIMO Equalization'' section stating that three
equalizers (ZF, LMMSE, LMMSE-SIC) ``were implemented,'' with speed-up figures,
although MIMO is future work and no MIMO result appears anywhere in the thesis, and although the summary of removed material in Appendix~E already listed the
$2\times2$ MIMO second hop and its equalization among the material deleted in the
canonical restructure. The section had simply survived that deletion.
\emph{Resolved:} the section is now \emph{removed}, matching what Appendix~E
already recorded. The equalizer implementations themselves have since been deleted
from the framework as well, so no unused MIMO code remains (Appendix~C), and with the
section gone the equalizer-terminology question raised by the supervisor
(Appendix~E, Chapter~2 comment~7) no longer arises: ``MMSE'' now appears in this
thesis only in its correct non-linear sense, denoting the MMSE-optimal estimator
$\mathbb{E}[x\mid y]$ of the relay denoising task.

\end{enumerate}
\endgroup

\subsection*{F.2 Consistency corrections}

\begingroup\small
\begin{enumerate}\setlength{\itemsep}{3pt}\setlength{\parskip}{1pt}
\setcounter{enumi}{11}

\item \textbf{Missing English abstract.} \texttt{main.tex} included a nonexistent
file, so the abstract never reached the compiled PDF. \emph{Resolved:} the include
now points at \texttt{chapters/frontmatter.tex}.

\item \textbf{Monte-Carlo budget stated as uniform.} The abstract claimed one
budget for all results, but E1 uses $20\times50$k, E2--E3 $10\times10$k, and the
unknown-channel study $5\times50$k. \emph{Resolved:} per-experiment budgets are
tabulated in Section~\ref{sec:monte-carlo-ber-estimation}, the abstract is scoped
accordingly, and the caveat that a $5$-trial Wilcoxon test cannot reach $p<0.05$
is stated.

\item \textbf{Discussion contradicted Table~\ref{tbl:table2}.} The low-SNR
architecture ordering claimed the sequence models won, whereas the canonical table
shows the minimal feedforward relays leading. \emph{Resolved:} the paragraph now
matches the table, which also strengthens the H3/H4 narrative.

\item \textbf{Training-time and unit drift.} Training times were quoted from a
different run than Table~\ref{tbl:table13}; a normalized-comparison gap appeared
as ``$\sim$1~dB'' in one chapter and ``$\sim$1\% BER'' in another; the MLP
training time was rounded to ``under 3~seconds'' against a tabulated 4.9~s.
\emph{Resolved:} Table~\ref{tbl:table13} is the single source of truth throughout,
and the units and roundings now agree across chapters.

\item \textbf{Orphaned results and conclusions.} A CSI/LayerNorm table and three
conclusions referenced a 48-variant search, 16-PSK results, and a nonexistent
section, all removed in the restructure. \emph{Resolved:} the table and the
three conclusions were deleted (see Appendix~E).

\item \textbf{Figure and table disagreed with the committed data.} The 16-class
grouped-bar figure came from a superseded run, and several cells of
Table~\ref{tbl:table24} traced to the same stale source. \emph{Resolved:} the
figure was regenerated from the committed results file and every cell of the table
was reset from it, so table, figure, and caption now share one source.

\item \textbf{Bibliography defects.} One entry carried an incorrect author, title,
and page range; two works appeared twice under different keys; and the Mamba
architecture was attributed to the earlier S4 paper. \emph{Resolved:} the entry was
corrected, the duplicates merged, and the citation repointed.

\item \textbf{Imprecise figures in the summary.} The summary said ``a few dB
behind a pilot-aided Viterbi receiver'' where the body reports $\approx$1--1.5~dB
(pure ISI) and $\approx$2~dB (composite), and a reported confidence interval was
quoted at the wrong SNR point. \emph{Resolved:} both are now stated precisely and
consistently with the underlying results files.

\end{enumerate}
\endgroup

\subsection*{F.3 Editorial corrections}

\begingroup\small
\begin{enumerate}\setlength{\itemsep}{3pt}\setlength{\parskip}{1pt}
\setcounter{enumi}{19}

\item \textbf{Misdirected equation pointers.} Several parenthetical
``(Equation~$\ldots$)'' pointers referenced the wrong equation, notably the
supervised-loss and $\beta$-VAE sentences. \emph{Resolved:} each now points at its
intended target, or the pointer was removed where it served no purpose.

\item \textbf{Noise-convention note could not be cross-referenced.} The note was a
\texttt{quote} environment whose label resolved to the enclosing section number.
\emph{Resolved:} it is promoted to a numbered remark
(Remark~\ref{noise-convension}) and both citing sites reference it properly.

\item \textbf{Malformed markup in the Mamba block.} A bare \texttt{\textbackslash
ref}, a missing space, an \texttt{\textbackslash addcontentsline} placed inside an
equation environment, and a duplicated entry. \emph{Resolved:} the block was
rewritten with correct structure and complete sentences.

\item \textbf{Broken enumeration and truncated caption.} The unknown-channel setup
list ran (i), (ii), (iv), (v); a figure caption ended mid-sentence; stray
characters and a double-encoded ``$\times$'' appeared in places; and one section
label still read ``partially-confirmed'' after the finding was confirmed.
\emph{Resolved:} all corrected.

\item \textbf{Summary item too dense.} The unknown-channel summary item ran to
roughly 250 words and buried the pilot-crossover and complexity results.
\emph{Resolved:} it is split into three items mirroring the structure of
Chapter~\ref{sec:unknown-channels}.

\end{enumerate}
\endgroup

\noindent\textbf{Findings recorded but requiring no change.} The review also
recorded points of agreement that entailed no edit: the symbol-wise-DF and
window-causality remarks (Remarks~\ref{rem:df-terminology}
and~\ref{rem:window-causality}) were noted as correctly pre-empting the obvious
objections; the analytic $0.25$ error floor and its four slicer amplitudes were
independently re-derived and confirmed; the flat-channel control was identified as
the methodological highlight of the unknown-channel chapter; and the bounded
phrasing of the H6 conclusions, never claiming superiority over a matched
classical receiver, was endorsed as the defensible reading of the evidence.
